% Table for ch8-19 -- Conflict Resolution Across Harbors (S5.2), a
% classification (rule x defends against x why rejected/adopted).
%
% Reader question: why is "only the issuer may revoke its own cards" the
%   right default, and what do the alternatives get wrong?
% Claim: of three candidate rules, only the authoritative rule avoids both
%   failure modes (an attacker forging revocations, and honest cards
%   rejected in equilibrium); the third rule is a degraded mode of the
%   second, not an independent option.
% Evidence: S:fh-revocation-conflict's own three-item enumeration.
% Grammar rejected: a decision-tree figure would suggest the three rules
%   are evaluated sequentially; they are evaluated in parallel against the
%   same attacker, which a table's independent rows states directly.
% Five-second test: which rule is adopted, and which one does it subsume
%   as a degraded mode?
\begin{table}[H]
\centering
\caption{Three candidate revocation-conflict rules, evaluated against the
same disagreement (harbor $A$ has revoked card $c$; harbor $B$ has not yet
seen it). The adopted rule is the only one with a bounded attacker window
under stated assumptions. [internal]}
\label{tab:fh-revocation-rules}
\small
\renewcommand{\arraystretch}{1.2}
\begin{tabularx}{\linewidth}{@{}>{\raggedright\arraybackslash}p{0.24\linewidth} >{\raggedright\arraybackslash}X >{\raggedright\arraybackslash}X@{}}
\toprule
\textbf{rule} & \textbf{what it defends against} & \textbf{why rejected / adopted} \\
\midrule
pessimistic (anyone revokes) & an issuer that refuses to revoke & rejects too many honest cards; an attacker injects spurious revocations into gossip \\
\textbf{authoritative (only issuer revokes)} & \textbf{spurious revocations from a non-issuer} & \textbf{adopted: an \mbox{operator-priced} attacker window under stated \mbox{service-level} terms} \\
epoch-bounded ($\Delta$-epoch deadline) & a harbor persistently falling behind & a degraded mode of authoritative with a watchdog, not a separate rule \\
\bottomrule
\end{tabularx}
\end{table}
